\subsection{Runtime Enforcement and Invariant Guards}

\noindent \textbf{Invariant Generation and Enforcement.}
Prior research has proposed multiple approaches to derive and enforce smart-contract invariants.
Cider\cite{liu2022learning} learns arithmetic-overflow invariants using reinforcement learning,
InvCon and InvCon+\cite{liu2022invcon,liu2024automated} combine dynamic inference with static validation,
and Trace2Inv\cite{chen2024demystifying} studies transaction-history-based invariant learning.
These directions provide practical invariant families and evidence for runtime protection,
especially when the objective is to reject abnormal executions with low deployment friction.

\noindent \textbf{Control-Flow Restriction in Practice.}
Industry systems such as SphereX use control-flow restriction to harden deployed protocols\cite{spherex}.
This line of work emphasizes operationally actionable policy gates,
including whitelist/blacklist style control decisions over execution paths.

\noindent \textbf{Re-entrancy Defense and Type-System Approaches.}
Many tools target re-entrancy-specific control-flow risks using static and symbolic analyses,
including Oyente, Osiris, ReGuard, Slither, MPro, Sailfish, Pluto, Park, SliSE,
and callback-oriented analyses\cite{luu2016making,torres2018osiris,liu2018reguard,feist2019slither,zhang2019mpro,bose2022sailfish,ma2021pluto,zheng2022park,wang2024efficiently,albert2020taming}.
Runtime defenses such as Sereum and online callback monitoring also protect deployed contracts\cite{rodler2018sereum,grossman2017online},
while temporary interaction restrictions are explored for reducing dangerous call patterns\cite{callens2024temporarily}.
In parallel, secure type-system research formalizes information-flow and compositional re-entrancy guarantees\cite{cecchetti2020securing,cecchetti2021compositional}.
Taken together, these works motivate broader runtime policy frameworks that can go beyond single-vulnerability families.
